% !TeX root = ../navier-stokes-manuscript.tex
% ============================================================
% 2.2 Finite-Time Response Escalation
% ============================================================

\subsection{Finite-Time Response Escalation}\label{sub:FiniteTimeResponseEscalation}

Set
\[
  V_j(t)
  :=
  \partial_t^j u(\,\cdot\,,t),
  \qquad
  j\in\Nzero.
\]
These are Eulerian time derivatives: the spatial coordinate is held fixed while time changes.
Let \(X\) be a Banach space, fix \(0\leq t_0<T_*\), and suppose each row under consideration is absolutely continuous as an \(X\)-valued map on every interval \([t_0,t]\) with \(t<T_*\), with \(V_{j+1}\) the Bochner derivative of \(V_j\).
For any fixed \(j\), the fundamental theorem of calculus gives
\[
  V_j(t)
  =
  V_j(t_0)
  +
  \int_{t_0}^{t}V_{j+1}(\tau)\,d\tau.
\]
Taking the \(X\)-norm gives
\[
  \norm{V_j(t)}_X
  \leq
  \norm{V_j(t_0)}_X
  +
  \int_{t_0}^{t}\norm{V_{j+1}(\tau)}_X\,d\tau.
\]
Thus
\[
  \int_{t_0}^{T_*}\norm{V_{j+1}(\tau)}_X\,d\tau
  <
  \infty
\]
gives a finite upper bound for \(V_j\) throughout the remaining interval.
It also gives a unique endpoint value in \(X\), since for \(t_0<s<t<T_*\),
\[
  \norm{V_j(t)-V_j(s)}_X
  \leq
  \int_s^t\norm{V_{j+1}(\tau)}_X\,d\tau.
\]
The absolute continuity of the integral makes \(V_j(t)\) Cauchy as \(t\uparrow T_*\).

Now suppose that the \(j\)-th row becomes unbounded in \(X\):
\[
  \sup_{t_0\leq t<T_*}
  \norm{V_j(t)}_X
  =
  \infty.
\]
The preceding estimate forces
\[
  \int_{t_0}^{T_*}
  \norm{V_{j+1}(\tau)}_X
  \,d\tau
  =
  \infty.
\]
Since \(T_*-t_0\) is finite,
\[
  \int_{t_0}^{T_*}
  \norm{V_{j+1}(\tau)}_X
  \,d\tau
  \leq
  (T_*-t_0)
  \sup_{t_0\leq t<T_*}
  \norm{V_{j+1}(t)}_X
\]
whenever the supremum on the right is finite.
Consequently,
\[
  \sup_{t_0\leq t<T_*}
  \norm{V_j(t)}_X
  =
  \infty
  \quad\Longrightarrow\quad
  \sup_{t_0\leq t<T_*}
  \norm{V_{j+1}(t)}_X
  =
  \infty.
\]
Repeating the same argument gives
\[
  \sup_{t_0\leq t<T_*}
  \norm{V_{j+k}(t)}_X
  =
  \infty
  \qquad
  \text{for every finite }k\geq1.
\]

A finite-time blow-up at one rung must be accompanied by unbounded responses at every rung above it.
The statement begins at whichever Eulerian row first becomes unbounded in the chosen norm.
Applying the same calculation to \(\partial_x^\alpha V_j\) shows that a gradient or any higher mixed coordinate may be the first escaping quantity; the escalation then begins at that coordinate.
The implication concerns size alone.
Each higher row may oscillate, change direction, and pass through zero while still having the unbounded size required by the implication.

The physical language of acceleration and jerk follows a moving fluid particle, whereas \(V_1\) and \(V_2\) hold space fixed.
The distinction matters because transport separates these two kinds of differentiation.
The original equation tells us exactly how they fit together, and differentiating that equation produces the full response family forced by the fluid itself.
